\documentclass[11pt]{article}

\usepackage[margin=1in]{geometry}
\usepackage{amsmath}
\usepackage{amssymb}
\usepackage{mathtools}
\usepackage{xcolor}
\usepackage[colorlinks=true,linkcolor=blue,urlcolor=blue]{hyperref}
\usepackage[final]{microtype}

\title{heterosplat: Viewer Architecture}
\author{Source: \texttt{repos/heterosplat/Documents/LaTeX/ViewerArchitecture.tex}}
\date{\today}

\newcommand{\R}{\mathbb{R}}

\begin{document}
\maketitle

\section{Purpose}

This document describes the architecture and rendering pipeline of the
\textbf{SplatViewer} (Phase~3 of heterosplat). The viewer provides an
interactive GLFW + Dear~ImGui interface for viewing and transforming
3D Gaussian splats at interactive frame rates.

\paragraph{Source files.}
\begin{itemize}
  \item \texttt{Source/Viewer/SplatViewer.cu} --- main viewer binary (GLFW
        window, ImGui panels, render loop, camera controls).
  \item \texttt{Source/Viewer/ColormapDebug.h} / \texttt{ColormapDebug.cu} ---
        second custom CUDA kernel: per-axis colormap for debug visualization.
  \item \texttt{Source/Viewer/CMakeLists.txt} --- builds ImGui (FetchContent)
        and ColormapDebug library.
\end{itemize}

\section{Rendering pipeline per frame}

Each frame executes the full forward rasterization pipeline on the GPU:

\begin{enumerate}
  \item \textbf{Transform.} Copy original Gaussian parameters to working
        buffers. Apply the user's similarity transform (rotation +
        scale + translation) via \texttt{apply\_similarity\_transform\_kernel}
        from the Normalize module.

  \item \textbf{Activate.} Convert log-scales to actual scales via
        \texttt{exp\_forward} (opacities are activated once at load time
        since they don't change with the transform widget).

  \item \textbf{View directions.} Compute per-Gaussian view directions
        from camera position to mean positions.

  \item \textbf{Projection.} Run EWA 3DGS fused projection to compute
        2D means, depths, conics, and bounding radii.

  \item \textbf{Spherical harmonics.} Evaluate SH coefficients with view
        directions to produce per-Gaussian RGB colors.

  \item \textbf{Tile intersection (two-pass).}
        Pass~1: count intersections per tile.
        CUB inclusive prefix sum to get cumulative tile counts.
        Pass~2: write intersection IDs and flatten IDs.
        CUB radix sort by tile+depth.

  \item \textbf{Tile offsets.} Compute per-tile start offsets into the
        sorted intersection array.

  \item \textbf{Rasterize.} Alpha-composite sorted Gaussians per tile
        to produce the final RGB image.

  \item \textbf{Display.} Download the CUDA render buffer to host memory,
        upload to an OpenGL texture via \texttt{glTexSubImage2D}, and draw
        a fullscreen quad.
\end{enumerate}

\section{Camera model}

The viewer uses an orbit camera parameterized by:
\begin{itemize}
  \item \textbf{Center point} $(c_x, c_y, c_z)$ --- the look-at target.
  \item \textbf{Azimuth} $\theta$ --- horizontal angle (degrees).
  \item \textbf{Elevation} $\phi$ --- vertical angle (degrees), clamped to
        $(-89°, 89°)$.
  \item \textbf{Radius} $r$ --- distance from center.
  \item \textbf{FOV} --- vertical field of view (degrees).
\end{itemize}

Camera position in world space:
\begin{equation}
  \mathbf{p}_\text{cam} = \mathbf{c} +
  r \begin{pmatrix}
    \cos\phi\,\cos\theta \\
    \cos\phi\,\sin\theta \\
    \sin\phi
  \end{pmatrix}.
\end{equation}

The view matrix is constructed using a Z-up look-at convention
(matching COLMAP) with the camera's Y-axis pointing downward in screen
space.

\paragraph{Controls.}
\begin{itemize}
  \item Left-drag: orbit (azimuth + elevation).
  \item Right-drag: pan (translate center).
  \item Scroll wheel: zoom (scale radius).
\end{itemize}

\section{Transform widget}

The ImGui ``Transform'' panel exposes per-source:
\begin{itemize}
  \item Translation $\mathbf{t} \in \R^3$ (sliders, range $[-5, 5]$).
  \item Rotation $(\alpha, \beta, \gamma)$ in degrees, applied as
        $R = R_z(\gamma)\,R_y(\beta)\,R_x(\alpha)$.
  \item Uniform scale $s$ (slider, range $[0.1, 5.0]$).
\end{itemize}

A ``Reset'' button returns all transform parameters to identity.

\section{Colormap debug kernel}
\label{sec:colormap}

The \texttt{colormap\_per\_axis\_kernel} is the second custom CUDA kernel
in heterosplat (the first being \texttt{apply\_similarity\_transform\_kernel}
from Phase~2).

Given a selected coordinate axis $a \in \{X, Y, Z\}$ and the scene's
bounding box $[\min_a, \max_a]$, the kernel maps each Gaussian's position
along that axis to a normalized parameter $t \in [0, 1]$:
\begin{equation}
  t_n = \frac{x_{n,a} - \min_a}{\max_a - \min_a}.
\end{equation}

This parameter is then mapped through a piecewise-linear approximation
of the Turbo colormap, producing an RGB color $(r, g, b)$. The color is
written to the Gaussian's SH DC coefficient as:
\begin{equation}
  c_{\text{DC}} = \frac{(r, g, b)}{C_0},
  \qquad
  C_0 = \frac{1}{2\sqrt{\pi}} \approx 0.28209.
\end{equation}

Higher-order SH coefficients ($\ell \geq 1$) are zeroed so that the
colormap is view-independent.

\paragraph{Use case.} Toggle the ``Colormap by axis'' checkbox in the
Debug panel to visualize how Gaussians are distributed along each
coordinate axis. This is useful for verifying coordinate convention
alignment when fusing multiple sources.

\section{Dependencies}

\begin{center}
\begin{tabular}{ll}
  \textbf{Library} & \textbf{Version / source} \\
  \hline
  GLFW & 3.3.6 (system \texttt{libglfw3-dev}) \\
  GLEW & 2.2.0 (system \texttt{libglew-dev}) \\
  Dear ImGui & v1.91.8 (CMake FetchContent) \\
  OpenGL & Core profile 3.3 \\
\end{tabular}
\end{center}

\end{document}
